\section{Appendix}
\label{sec:appendix}


% \subsection{TBD}

% \xhdr{Latent overthinking}
% To analyze latent thinking patterns, we verbalize tokens from all iteration depths using their last-layer hidden states.
% The oracle method uses the oracle policy $\pi$ from Section~\ref{sec:training} for iteration decision.
% (1) \textbf{Generation.} 
% Since ground-truth tokens are unavailable during generation, we use predictions from the stronger DeepSeek-R1-Distill-Qwen-32B model~\citep{guo2025deepseek-r1} as proxy labels.
% Table~\ref{tab:oracle} shows that the oracle policy substantially improves performance by verbalizing correct predictions immediately while iterating only on incorrect ones. 
% With our trained iteration decider approximating the oracle, \name outperforms both Standard and AlwaysThink baselines. 
% However, the ideal oracle policy achieves even higher gains, indicating future potential.
% (2) \textbf{Next-token prediction}. 
% We evaluate next-token prediction accuracy on the Open-R1 dataset, using the test model itself as the reference model in $\pi$.
% Figure~\ref{fig:idea} reveals that AlwaysThink produces more incorrect than correct revisions, demonstrating latent overthinking.
% In contrast, oracle-controlled iterations substantially increase correct revisions by selectively targeting hard tokens.

% Figure
% Objective: Performance w.r.t. iter count
% X: iter count (threshold)
% Y: benchmark accuracy (choose the best performing one)

%%% continuation threshold and flow (separate table and figure) %%%





%%%%%%%%%%%% Experiment setup %%%%%%%%%%%%
\subsection{Additional Experiment Setups}
\label{sec:appendix/exp_setup}

\subsubsection{Training Recipe}
\label{sec:appendix/exp_setup/training_recipe}

\begin{wraptable}{r}{0.40\linewidth}
\vspace{-\baselineskip}
\caption{Training hyperparameters.}
    \label{tab:train_hparams}
    \centering
    \begin{tabular}{ll}
    \toprule
    \textbf{Hyperparameter} & \textbf{Value} \\
    \midrule
    learning rate & 4e-5 \\
    max grad norm & 0.2 \\
    training epochs & 5 \\
    global batch size & 128 \\
    warmup ratio & 0.03 \\
    lr scheduler & cosine (min-lr ratio 0.1) \\
    precision & bfloat16 \\
    \bottomrule
    \end{tabular}
\end{wraptable}

We follow the official training setup of Open-R1~\citep{openr1} and use the Mixture-of-Thoughts dataset.
For our main performance experiments (Section~\ref{sec:experiment/performance}), we filter samples with output length exceeding 8K tokens from each category (math, code, and QA), then randomly sample 33K examples from each to form a balanced 100K training set.
The filtered dataset contains 480M tokens, with 1\% reserved for validation. 
For design choice exploration (Section~\ref{sec:exp/dse}), we use only the math subset filtered by 8K output length, resulting in 75K training samples.
For 0.6B and 1.7B models, we use a maximum sequence length of 8192 tokens for all methods.
For 4B models, we reduce the maximum sequence length to 4096 for Standard, \name, TaH+, and 3072 for Ouro due to memory constraints.
Detailed training hyperparameters are listed in Table~\ref{tab:train_hparams}.

\subsubsection{Baseline Setups}
\label{sec:appendix/setup/baseline}

\xhdr{AlwaysThink}
AlwaysThink uses the same architecture as TaH, except that it iterates twice at every token position and uses standard causal attention (attending only to the current iteration depth) instead of duo-causal attention.

\xhdr{Ouro}
We implement the Ouro architecture following its official design~\citep{zhu2025ouro}, but fine-tune it on the Open-R1 dataset from Qwen-Base initialization, to align with all other methods.
During fine-tuning, we adopt the entropy-regularized loss of Ouro.
Since the original Ouro is trained from scratch on different data distributions, performance may differ from the original paper.

\subsubsection{Iteration Depth Flexibility}
\label{sec:appendix/oracle_policy}

The main experiments in this paper focus on $D_{\max}=2$.
Here we detail how the oracle iteration policy generalizes to arbitrary maximum iteration depth, as used in Section~\ref{sec:exp/dse}.
The focus of this work is how to \emph{selectively attribute} iterations to different tokens, rather than increasing the average iteration depth.

\xhdr{Binary case ($D_{\max}=2$)}
Prior work has proposed many ways to estimate token difficulty, including excess loss~\citep{lin2024rho, xie2023doremi}, entropy~\citep{wang2025beyond, chen2023eeLLM}, and prediction difference~\citep{fu2025r2r}.
For shallow iteration budgets of up to two, we adopt the prediction-difference policy used in the main experiments: tokens that fail top-1 next-token prediction at the first iteration are labeled as hard.
Formally, the oracle iteration depth $d_i^\pi$ follows a binary rule:
\begin{equation}
    d_i^\pi \;=\; \begin{cases}
        1, & \text{if } h_i = 0 \quad \text{(easy token)} \\
        D_{\max}, & \text{if } h_i = 1 \quad \text{(hard token)} ,
    \end{cases}
\end{equation}
where $h_i$ is the hard-token indicator.
This induces a Bernoulli distribution over depths: easy tokens verbalize at depth~1, hard tokens iterate to depth~$D_{\max}$.

\xhdr{General framework ($D_{\max}>2$)}
For deeper iteration budgets, we replace the binary rule with a continuous difficulty score and a target depth distribution.
We use the reference model's first-pass cross-entropy as a continuous measure of token difficulty:
\[
\ell_i^{\text{ref}} \;=\; -\log p_{i,\text{ref}}^{(1)}\!\big(t_{i+1}\big).
\]
Let $\Lambda$ be a target distribution over $\{1,\dots,D_{\max}\}$ specifying the desired fraction of tokens at each depth.
We assign iteration depths via quantile mapping: rank all tokens by $\ell^{\text{ref}}$ to obtain $u_i = \mathrm{QuantileRank}(\ell_i^{\text{ref}}) \in [0,1]$, then set
\begin{equation}
    d_i^\pi \;=\; F_\Lambda^{-1}(u_i),
\end{equation}
where $F_\Lambda^{-1}$ is the quantile function (inverse CDF) of $\Lambda$.
Harder tokens (higher loss) thus receive deeper iterations, and the overall depth distribution follows $\Lambda$.
The binary case is recovered when $\Lambda$ matches the depth distribution induced by the hard-token labels (i.e., mass at depths~1 and~$D_{\max}$ equal to the easy and hard fractions, respectively).
The oracle depth $d_i^\pi$ is then converted into per-depth supervision for the iteration decider (Section~\ref{sec:training}): at each depth $d \in \{1,\ldots,D_{\max}-1\}$, the continuation label $c_i^{(d)} = \mathbb{1}[d < d_i^\pi]$ equals $1$ (continue iterating) when the token has not yet reached its assigned depth, and $0$ (verbalize) once $d = d_i^\pi$.
At inference, the decider's predicted continuation probabilities $\hat{c}_i^{(d)}$ are compared against the threshold $c_{\text{threshold}}$ to decide whether each token continues at each depth.

\xhdr{Deeper-depth experiment setup}
For TaH-3-Decreasing in Table~\ref{tab:max_interation_3}, we set the training target distribution over depths to $\Lambda=(0.88, 0.06, 0.06)$.
For TaH-4-Decreasing, we set $\Lambda=(0.88, 0.04, 0.04, 0.04)$.
That is, we keep the depth-1 fraction the same as TaH-2 and evenly allocate the remaining budget to deeper depths; the inference continuation threshold is set to $c_{\text{threshold}}=0.9$, which makes the realized continuation rate decrease with depth.
For comparison, TaH-3-Uniform uses $\Lambda=(0.34, 0.33, 0.33)$ during training and $c_{\text{threshold}}=0.5$ during inference to preserve a roughly uniform depth allocation.
The realized continuation ratios are reported in brackets in Table~\ref{tab:max_interation_3}.

\subsubsection{Parameter Breakdown}
\label{sec:appendix/exp_setup/parameter_breakdown}

Table~\ref{tab:param_composition} reports the parameter breakdown of the Standard, \name, and TaH+ methods.
To offset the additional parameters introduced by \name through LoRA and the iteration decider, we remove one layer from the LLM backbone, ensuring a fair comparison.
In practical deployments, we recommend TaH+, which adds only about 3\% additional parameters.

\begin{table}[b]
    \caption{Parameter breakdown of Standard, TaH, and TaH+. Counts are reported using M (million) and B (billion).}
    \label{tab:param_composition}
    \centering
    \begin{tabular}{ll|cccc}
    \toprule
    \textbf{Param.} & \textbf{Method} & \textbf{Backbone} & \textbf{LoRA} & \textbf{Iter. Decider} & \textbf{Total} \\
    \midrule
    \multirow{3}{*}{\textit{0.6B}} & Standard & 596M & -- & -- & 596M \\
    & \name & 580M & 10M & 5M & 595M \\
    & TaH+ & 596M & 10M & 5M & 611M \\
    \midrule
    \multirow{3}{*}{\textit{1.7B}} & Standard & 1.72B & -- & -- & 1.72B \\
    & \name & 1.67B & 17M & 18M & 1.71B \\
    & TaH+ & 1.72B & 17M & 18M & 1.76B \\
    \midrule
    \multirow{3}{*}{\textit{4B}} & Standard & 4.02B & -- & -- & 4.02B \\
    & \name & 3.92B & 32M & 70M & 4.02B \\
    & TaH+ & 4.02B & 33M & 70M & 4.12B \\
    \bottomrule
    \end{tabular}
\end{table}

% \subsubsection{Latent Overthinking Analysis Setup}
% \label{sec:appendix/exp_setup/oracle}

% We set up the supplementary generation-time oracle experiment in Figure~\ref{fig:generation_oracle} to estimate the performance upper bound of our method. The oracle employs the DeepSeek-R1-Distill-Qwen-32B model as a dynamic label generator, replacing the MLP-based iteration decider.
% During each iteration, we compare the token predictions from the label generator with those from the \name model.
% The model continues to iterate only when the top-1 predictions of these two models differ.
% Due to resource constraints and computational overhead, we evaluated the accuracy only on the first 100 samples from the MATH500 dataset, denoted as MATH100 throughout the paper.

%%%%%%%%%%%% Experiments %%%%%%%%%%%%
\subsection{Additional Experimental Results}

\subsubsection{Math-Only Training}
This section isolates the math-only setting: models are trained on the math subset of Open-R1 and evaluated on math benchmarks as well as out-of-domain tasks.

%%% Main table %%%

% Table
% Objective: Math Performance at Equation #Parameter
% Setup
% dataset:
%   Accuracy: MATH500, AMC23, GSM8K, TBD(AIME24; AIME25), Olpmpiadbench
%   Perplexity: Open-R1 test set (validation set)
% method:
% Ours; SFT-base; AlwaysMaxIter
% Others:
% soft-thinking; DiT (pause token)
\begin{table*}[t]
    \caption{Accuracy comparison of different baselines across five benchmarks and three model sizes.
    Subscripts indicate improvement over Standard. The top two scores for each task and model size are highlighted in bold.
    }
    \label{tab:math_performance}
    \centering
    \setlength\tabcolsep{3pt} % default value: 6pt
    \begin{tabular}{ll|cccccc}
    \toprule
    & & \multicolumn{6}{c}{\textbf{Method}} \\
    \cmidrule(lr){3-8}
    \textbf{Param.} & \textbf{Benchmark} & \textbf{Standard} & \textbf{Routing}& \textbf{SoftThink}& \textbf{AlwaysThink}& \textbf{\name} & \textbf{TaH+} \\
    \midrule
    \multirow{7}{*}{\textit{0.6B}}
    & AIME25 & \textbf{4.2} & 1.0 & 2.5 & 1.5 & \textbf{4.2} & \textbf{5.0} \\
    & OlympiadBench & 18.8 & 7.4  & 19.4 & 10.2 & \textbf{23.9}  & \textbf{24.0} \\
    & AMC23 & 23.4 & 10.9 & 24.1 & 15.6 & \textbf{32.5} &  \textbf{30.6} \\
    & MATH500 & 47.2 & 27.3 & 48.8 & 32.8 & \textbf{51.2} & \textbf{54.2} \\
    & GSM8K & 62.5 & 45.6  & 61.3 & 54.6 & \textbf{64.4} & \textbf{68.8} \\
    % & MMLU-STEM & 51.6 & 37.9 & 51.4 & 42.6 & 50.8 & \textbf{56.3} \\
    % & AIME24 & 2.7 &  & 0.8 &  & \textbf{3.3} & \textbf{2.1} \\
    % & Minerva & 9.2 &  &  &  & \textbf{10.3} & \textbf{13.2} \\
    \cmidrule(lr){2-8}
    & Average & 31.2 & 18.5 & 31.2 & 22.9 & \textbf{35.2{\scriptsize$/+$4.0}} & \textbf{36.5{\scriptsize$/+$5.3}} \\
    % & Avg. improvement &  & & & & & \\

    \midrule
    \multirow{7}{*}{\textit{1.7B}}& AIME25 & 13.3 & 10.2 & 12.9 & 10.0 & \textbf{17.9} & \textbf{14.6} \\
    & OlympiadBench & 33.0 & 30.6 & 33.4 & 30.0 & \textbf{38.8} & \textbf{41.2} \\
    & AMC23 & 42.2 & 42.2 & 43.1 & 42.5 & \textbf{48.4} & \textbf{51.2} \\
    & MATH500 & 68.4 & 60.0 & 68.8 & 61.8 & \textbf{74.4} & \textbf{73.0} \\
    & GSM8K & 82.1 & 71.2 & 79.6 & 79.3 & \textbf{84.5} & \textbf{85.8} \\
    % & MMLU-STEM & 70.8 & 61.7 & 70.6 & 63.8 & 69.9 & \textbf{73.7} \\
    % & AIME24 & 9.4 &  & 10.0 &  & \textbf{10.8} & \textbf{12.1} \\
    % & Minerva & 25.4 &  &  &  & \textbf{27.2} & \textbf{29.4} \\
    \cmidrule(lr){2-8}
    & Average & 47.8  & 36.8 & 47.6 & 44.7 & \textbf{52.8{\scriptsize$/+$5.0}}& \textbf{53.2{\scriptsize$/+$5.4}}\\

    \midrule
    \multirow{7}{*}{\textit{4B}}& AIME25 & 23.3& 22.5 & 22.5 & \multirow{5}{*}{OOM}& \textbf{30.4}& \textbf{28.3} \\
    & OlympiadBench & 47.7& 45.0 & 50.1 & & \textbf{50.5}& \textbf{52.0} \\
    & AMC23 & 62.8& 60.9 & 64.1 & & \textbf{70.3}& \textbf{70.6} \\
    & MATH500 & 82.8& 76.1 & 83.2 & & \textbf{84.4}& \textbf{85.6} \\
    & GSM8K & 90.5& 85.3 & \textbf{90.9} & & 90.4& \textbf{91.5} \\
    % & MMLU-STEM & 84.1 &  & \textbf{84.4} &  & 83.3 & 84.4 \textbf{} \\
    % & AIME24 & 9.4 &  & 10.0 &  & \textbf{10.8} & \textbf{12.1} \\
    % & Minerva & 25.4 &  &  &  & \textbf{27.2} & \textbf{29.4} \\
    \cmidrule(lr){2-8}
    & Average & 61.4& 58.0 & 62.2& --& \textbf{65.2{\scriptsize$/+$3.8}}& \textbf{65.6{\scriptsize$/+$4.2}} \\
    \bottomrule
    \end{tabular}
\end{table*}

\xhdr{Math benchmark evaluation}
In this section, all models are trained on the math subset of Open-R1 and evaluated on math benchmarks (see Section~\ref{sec:appendix/exp_setup/training_recipe} for training details).
We also add a \emph{Routing} baseline, which selects a model from a candidate pair for each question. In our experiments, we use two pairs: (1) MobileLLM-R1-360M, Qwen3-1.7B, and (2) Qwen3-0.6B, Qwen3-4B.
All candidate models are SFT-trained under the same settings. For each pair, the routing ratio is calibrated so that the average active parameter count matches our 0.6B and 1.7B targets, respectively.
Table~\ref{tab:math_performance} reports results for 0.6B, 1.7B, and 4B backbones across five challenging math benchmarks.
Even with strong Qwen3-Base initialization, existing approaches show limited effectiveness: AlwaysThink and routing methods fail to consistently outperform the standard baseline, while SoftThink yields only marginal gains.
In contrast, \name achieves stable improvements, with average gains of 4.0\% (0.6B), 5.0\% (1.7B), and 3.8\% (4B) over Standard.
TaH+, which adds less than 3\% more parameters, further improves to 5.3\%, 5.4\%, and 4.2\%, respectively.
For 0.6B and 1.7B, \name achieves 8.1-11.3\% gains over AlwaysThink, and TaH+ achieves 8.5-12.6\% gains. AlwaysThink-4B is not evaluated due to out-of-memory during training.

\xhdr{Out-of-domain evaluation}
We further evaluated the zero-shot generalization capability of models trained solely on math datasets from the main paper. 
As shown in Table~\ref{tab:math_1.7_peformance}, TaH+ demonstrates consistent improvements not only on in-domain math benchmarks (MATH500, AMC23) but also on out-of-domain tasks like MMLU-STEM. This indicates that the thinking patterns learned by TaH+ on math problems are robust and transferrable to broader scientific reasoning tasks.

% \subsubsection{Generalizability}
% \label{sec:appendix/exp/generalization}

% \xhdr{General training and evaluation}
% To verify generalizability, we expanded our training and evaluation to diverse domains. We followed the exact protocol from the main paper to fine-tune Qwen3-1.7B-Base. The only modification was replacing the math-only dataset with a balanced subset of OpenR1 (100k samples total) covering all task splits, ensuring a fair comparison by maintaining the same data scale.
% As shown in Table~\ref{tab:general_1.7_peformance}, TaH+ achieves consistent improvements across math, QA, and coding domains, with an average performance gain of 6.8\%. This demonstrates \name's effectiveness on diverse reasoning and generation tasks beyond pure mathematics.


% To assess cross-domain effectiveness, we train a 1.7B TaH model on the science subset of Open-R1. We exclude responses longer than 4,096 tokens, yielding 223M tokens in total; all other training settings follow Section~\ref{sec:experiment}. We then evaluate TaH against the Standard baseline on the GPQA-diamond benchmark~\citep{rein2023gpqa}. TaH improves performance from \textbf{35.4} to \textbf{39.9}, demonstrating the generalizability of our method.

% \begin{table}[tb]
%     \centering
%     \caption{Performance of Qwen3-1.7B models trained on a general OpenR1 mixture (math, QA, and code) across downstream benchmarks.}
%     \label{tab:general_1.7_peformance}
%     \begin{tabular}{ll|cccc}
%     \toprule
%     \textbf{Category} & \textbf{Dataset} & \textbf{Standard} & \textbf{SoftThink} & \textbf{AlwaysThink} & \textbf{TaH+} \\
%     \midrule
%     \multirow{2}{*}{Math} 
%         & MATH500 & 67.8 & 64.8 & 63.2 & \textbf{72.6} \\
%         & AMC23 & 39.7 & 40.3 & 40.9 & \textbf{48.4} \\
%     \midrule
%     \multirow{2}{*}{QA}
%         & GPQA & 30.3 & 33.3 & 30.5 & \textbf{39.4} \\
%         & MMLU-STEM & 74.1 & 73.5 & 69.6 & \textbf{76.6} \\
%     \midrule
%     \multirow{2}{*}{Code}
%         & HumanEval+ & 44.2 & 44.5 & 25.6 & \textbf{48.2} \\
%         & MBPP+ & 27.2 & 27.8 & 16.4 & \textbf{39.0} \\
%     \midrule
%     \multicolumn{2}{l|}{Average} & 47.2 & 47.4 & 41.0 & \textbf{54.0{\scriptsize$/+$6.8}} \\
%     \bottomrule
%     \end{tabular}
% \end{table}




\begin{table}[tb]
    \centering
    \caption{Performance of math-only trained models (0.6B and 1.7B) on in-domain math benchmarks and the out-of-domain STEM benchmark (MMLU-STEM).}
    \label{tab:math_1.7_peformance}
    \setlength{\tabcolsep}{4pt}
    \begin{tabular}{ll|cccc}
    \toprule
    \textbf{Param.} & \textbf{Benchmark} & \textbf{Standard} & \textbf{SoftThink} & \textbf{AlwaysThink} & \textbf{TaH+} \\
    \midrule
    \multirow{4}{*}{\textit{0.6}} 
        & MATH500    & 47.2 & 48.8 & 32.8 & \textbf{54.2} \\
        & AMC23      & 23.4 & 24.1 & 15.6 & \textbf{30.6} \\
        & MMLU-STEM  & 51.6 & 51.4 & 42.6 & \textbf{56.3} \\
        \cmidrule(lr){2-6}
        & Average    & 40.7 & 41.4 & 30.3 & \textbf{47.0} \\
    \midrule
    \multirow{4}{*}{\textit{1.7}} 
        & MATH500    & 68.4 & 68.8 & 61.8 & \textbf{73.0} \\
        & AMC23      & 42.2 & 43.1 & 42.5 & \textbf{51.2} \\
        & MMLU-STEM  & 70.8 & 70.6 & 63.8 & \textbf{73.7} \\
        \cmidrule(lr){2-6}
        & Average    & 60.5 & 60.8 & 56.0 & \textbf{66.0} \\
    \bottomrule
    \end{tabular}
\end{table}


\subsubsection{Additional Oracle Analysis}
\label{sec:appendix/generation_oracle}

We provide two complementary oracle analyses.
(1) \textbf{Downstream generation.}
Table~\ref{tab:generation_oracle_acc} reports generation results on MATH100 with Qwen3-0.6B.
We use DeepSeek-R1-Distill-Qwen-32B predictions~\citep{guo2025deepseek-r1} as proxy labels.
The oracle verbalizes when the model's top-1 prediction matches the proxy label, and iterates otherwise.
With our trained iteration decider approximating the oracle, \name outperforms both Standard and AlwaysThink baselines.
However, the ideal oracle policy achieves even higher gains, indicating future potential.
(2) \textbf{Next-token prediction.}
Figure~\ref{fig:generation_oracle_transition} shows prediction transitions obtained by verbalizing tokens from all iteration depths.
It reveals that AlwaysThink produces more incorrect than correct revisions, demonstrating latent overthinking.
In contrast, oracle-controlled iterations substantially increase correct revisions by selectively targeting hard tokens.

\begin{center}
    \begin{minipage}[t]{0.42\textwidth}
        \vspace{0pt}
        \captionsetup{hypcap=false}
        \captionof{table}{MATH100 accuracy under different training and inference policies on Qwen3-0.6B, using DeepSeek-R1-Distill-Qwen-32B predictions as proxy labels.}
        \centering
        \setlength\tabcolsep{4pt}
        \begin{tabular}{llc}
        \toprule
        \textbf{Training} & \textbf{Inference} & \textbf{Accuracy} \\
        \midrule
        Standard    & Standard    & 52 \\
        \midrule
        AlwaysThink & AlwaysThink & 38{\scriptsize$/-$14} \\
        AlwaysThink & TaH-Oracle       & 77{\scriptsize$/+$25} \\
        \midrule
        TaH-Oracle   & TaH-Decider   & 54{\scriptsize$/+$\phantom{0}2} \\
        TaH-Oracle   & TaH-Oracle       & 80{\scriptsize$/+$28} \\
        \bottomrule
        \end{tabular}
        \label{tab:generation_oracle_acc}
    \end{minipage}
    \hfill
    \begin{minipage}[t]{0.54\textwidth}
        \vspace{0pt}
        \captionsetup{hypcap=false}
        \captionof{figure}{Next-token prediction transitions across iteration depths on Qwen3-0.6B.}
        \centering
        \includegraphics[width=\linewidth]{figure/overthinking.pdf}
        \label{fig:generation_oracle_transition}
    \end{minipage}
\end{center}


\subsubsection{Additional Baselines}
\xhdr{Standard-pruned baseline}
Since \name removes one backbone layer to match the parameter budget of Standard, we additionally compare against a Standard-Pruned baseline with the same layer removed but without selective iteration.
Table~\ref{tab:standard_pruned} shows that pruning alone slightly degrades performance, while \name substantially outperforms both Standard and Standard-Pruned under the same training and evaluation setup.
This isolates the gain from selective latent iteration rather than from pruning.

\begin{table}[tb]
    \centering
    \caption{Comparison with a Standard-Pruned baseline on 1.7B math-trained models. Standard-Pruned removes the same backbone layer as \name but does not use latent iteration.}
    \label{tab:standard_pruned}
    \setlength{\tabcolsep}{5pt}
    \begin{tabular}{lcccccc}
    \toprule
    \textbf{Method} & \textbf{AIME25} & \textbf{Olympiad} & \textbf{AMC23} & \textbf{MATH500} & \textbf{GSM8K} & \textbf{Average} \\
    \midrule
    Standard & 13.3 & 33.0 & 42.2 & 68.4 & 82.1 & 47.8 \\
    Standard-Pruned & 11.7 & 32.7 & 41.3 & 68.0 & 79.4 & 46.6 \\
    \name & \textbf{17.9} & \textbf{38.8} & \textbf{48.4} & \textbf{74.4} & \textbf{84.5} & \textbf{52.8} \\
    \bottomrule
    \end{tabular}
\end{table}

\xhdr{Additional latent thinking baselines}
Some latent thinking methods require pre-training and use base models other than Qwen3.
We also compare with these methods, including Ponder~\citep{zeng2025pondering}. Specifically, we adopt the released pretrained PonderingPythia-1.4B as the base model and perform SFT on the same training data. We observe that the fine-tuned model learns the stylistic patterns of the training data,
but still underperforms substantially, which may be attributable to the limited capability of the PonderingPythia-1.4B backbone.

\begin{table}[tb]
\caption{Performance on MATH500 and GSM8K-500 (first 500 GSM8K samples)}
\label{tab:ponder_performance}
\centering
    \begin{tabular}{l|cc}
    \toprule
     & \multicolumn{2}{c}{\textbf{Method}} \\
     \cmidrule(lr){2-3}
     \textbf{Dataset} & Standard-0.6B & Ponder-1.4B \\
    \midrule
     MATH500 & 47.2 & 2.0 \\
     GSM8K-500 & 62.8 &  1.8  \\
    % & Minerva &  &   \\
    \midrule
     Avg. & 55 &  1.9   \\
    \bottomrule
    \end{tabular}
\end{table}

\subsubsection{Real-world Efficiency}
\label{sec:appendix/exp/real_world_efficiency}

\xhdr{Setup}
We investigate the real-world efficiency of different 1.7B models under our current implementation.
All measurements were obtained on a single A800 GPU with a batch size of 1 and a maximum output length of 8192 tokens, using a challenging AIME25 problem where all three methods reached the token limit.
Memory usage was profiled using \texttt{torch.cuda.memory.\_record\_memory\_history}.

\xhdr{Latency breakdown}
We report the decoding latency, throughput, and a detailed time breakdown for Standard, AlwaysThink, and \name in Table~\ref{tab:latency_breakdown}.
Here, \textit{Iter-1 forward} and \textit{Iter-2 forward} denote the total forward-pass time spent on the first and second latent iterations, respectively; \textit{Iter decider} is the time for the iteration decider network to judge whether to continue iterating or verbalize; \textit{LoRA switching} is the overhead of switching LoRA adapters; and \textit{Other} includes tensor initialization, concatenation, and related bookkeeping.

\begin{table}[t]
    \caption{Per-component latency breakdown on a single A800 GPU.}
    \label{tab:latency_breakdown}
    \centering
    \setlength\tabcolsep{4pt}
    \begin{tabular}{l|cc|cc|cc}
    \toprule
    \textbf{Component}
        & \multicolumn{2}{c|}{\textbf{Standard}}
        & \multicolumn{2}{c|}{\textbf{\name}}
        & \multicolumn{2}{c}{\textbf{AlwaysThink}} \\
        & \textbf{Latency (s)} & \textbf{Ratio(\%)} 
        & \textbf{Latency (s)} & \textbf{Ratio(\%)} 
        & \textbf{Latency (s)} & \textbf{Ratio(\%)} \\
    \midrule
    Iter-1 Forward & 210.6 & 100.0 & 229.8 & 76.2 & 224.1 & 30.0 \\
    Iter-2 Forward &  --   &  --   &  29.6 &  9.8 & 384.7 & 51.5 \\
    Iter. Decider   &  --   &  --   &  10.5 &  3.5 &  --   & --   \\
    LoRA Switching    &  --   &  --   &   7.5 &  2.5 &  91.1 & 12.2 \\
    Other          &  --   &  --   &  24.1 &  8.0 &  47.4 &  6.3 \\
    \bottomrule
    \end{tabular}
\end{table}

\xhdr{Discussion}
We note that our current implementation is not yet optimized at the system level, so there remains room for further efficiency improvements.
For example, the \textit{LoRA Switching} and \textit{Other} overheads (bookkeeping) are relatively high due to the Python-level implementation of dynamic control flow.
These engineering optimizations are largely orthogonal to the algorithmic design of \name, and we plan to continue refining the implementation to further reduce latency and memory overhead.
The theoretical FLOPs and memory access analysis of \name are provided in Appendix~\ref{sec:appendix/exp/theo_eff}.


\subsubsection{Theoretical Efficiency Analysis}
\label{sec:appendix/exp/theo_eff}

Following prior work~\cite{hoffmann2022training, yang2024glitches, wenheng2025cdllm}, we analyze the per-token decoding computation and memory access via analytical operator-level profiling: we trace the forward pass symbolically, accumulating FLOPs and memory bytes from tensor shapes and data types across all modules.

\xhdr{Notation and formulation}
Table~\ref{tab:eff_notation} lists the symbols used throughout this analysis;
Table~\ref{tab:op-cost} breaks down the cost from individual operators (\texttt{nn.linear}, \texttt{nn.sdpa} for \texttt{scaled\_dot\_product\_attention}) through per-layer modules to a full decode step. Since $h=n_h h_h$ for our LLMs, some terms are simplified accordingly.
Lightweight vector ops (RMSNorm, RoPE, SiLU, residual add) are also exhaustively profiled and included in the final numbers in Table~\ref{tab:efficiency_analysis_theo}. Since they generally contribute very little to the total cost, they are omitted from the formulas below for brevity.

\xhdr{Standard decoding}
For a standard LLM generating $t_\text{out}$ output tokens from a prefill of $t_\text{in}$ input tokens, the total FLOPs and memory access are:
\begin{align}
  \textstyle\sum_{i=1}^{t_\text{out}} \F_\text{step}(t_\text{in}{+}i{-}1), \qquad
  \textstyle\sum_{i=1}^{t_\text{out}} \M_\text{step}(t_\text{in}{+}i{-}1).
\end{align}
Per-token averages are obtained by dividing by $t_\text{out}$.

\begin{table}[ht]
\centering
\caption{Notation for theoretical efficiency analysis.}
\label{tab:eff_notation}
\footnotesize
\setlength{\tabcolsep}{4pt}
\begin{tabular}{@{}ll|ll@{}}
\toprule
\textbf{Symbol} & \textbf{Description} & \textbf{Symbol} & \textbf{Description} \\
\midrule
$h$   & Hidden dim              & $n_h$ & Query head count \\
$s$   & KV cache length         & $n_\text{kv}$ & KV head count \\
$d$   & Iteration depth         & $h_h$ & Per-head dim ($h/n_h$) \\
$l$   & Number of layers        & $h_\text{kv}$ & Total KV dim ($n_\text{kv} h_h$) \\
$w$   & Weight count per layer  & $h_\text{ff}$ & FFN intermediate size \\
$b$   & Bytes per element       & $t_\text{in}$ & Input token count \\
$v$   & Vocabulary size         & $t_\text{out}$ & Output token count \\
$k$   & Top-$k$ logits          & & \\
\bottomrule
\end{tabular}
\end{table}

\begin{table}[ht]
\centering
\caption{FLOPs and memory access per decode token, from operators to a full decode step. The Qwen architecture uses GQA ($g{=}n_h/n_\text{kv}$) and gated FFN (SwiGLU). Self-Attention applies 4 \texttt{nn.linear} (Q,K,V,O) plus \texttt{nn.sdpa}; FFN applies 3 \texttt{nn.linear} (gate, up, down). $w{=}2h(h{+}h_\text{kv}){+}3hh_\text{ff}$. Minor terms are absorbed into $\mathcal{O}(\cdot)$.}
\label{tab:op-cost}
\footnotesize
\setlength{\tabcolsep}{4pt}
\begin{tabular}{@{}lll@{}}
\toprule
\textbf{Component} & \textbf{FLOPs} & \textbf{Memory (bytes)} \\
\midrule
\multicolumn{3}{@{}l}{\textit{Operators}} \\[2pt]
\quad \texttt{nn.linear}
  & $2\,h_\text{in}\,h_\text{out}$
  & $(h_\text{in} h_\text{out} {+} h_\text{in} {+} h_\text{out})\,b$ \\[3pt]
\quad \texttt{nn.sdpa}
  & $4\,s\,h$
  & $[2h {+} s(2h_\text{kv} {+} 3n_h)]\,b$ \\
\midrule
\multicolumn{3}{@{}l}{\textit{Modules (per layer)}} \\[2pt]
\quad FFN
  & $6\,h\,h_\text{ff}$
  & $3\,h\,h_\text{ff}\,b + \mathcal{O}(hb)$ \\
\quad Self-Attention
  & $4h(h{+}h_\text{kv}) {+} 4sh$
  & $[2h(h{+}h_\text{kv}) {+} s(2h_\text{kv} {+} 3n_h)]\,b {+} \mathcal{O}(hb)$ \\[3pt]
\midrule
\multicolumn{3}{@{}l}{\textit{Decode step ($\times l$ layers, at length $s$)}} \\[2pt]
\quad Step
  & $l[4h(h{+}h_\text{kv}) {+} 4sh {+} 6hh_\text{ff}] {+} \mathcal{O}(lh_\text{ff})$
  & $l[w {+} s(2h_\text{kv} {+} 3n_h)]\,b {+} \mathcal{O}(lhb)$ \\
\bottomrule
\end{tabular}
\end{table}

\xhdr{Extension to \name}
\name assigns each of the $t_\text{out}$ output tokens ($i=1,\dots,t_\text{out}$) a depth $d_i\in\{1,\dots,d_\text{max}\}$.
Generating one output token at depth $d_i$ requires $d_i$ backbone passes, $d_i{-}1$ inner transitions (Equation~\ref{eq:ours_update}), and $\min(d_i,\,d_\text{max}{-}1)$ decider calls (one after each backbone pass; the call at $d_\text{max}$ is skipped because the token must verbalize).
The total decoding FLOPs and memory access over $t_\text{out}$ output tokens are:
\begin{align}
  &\sum_{i=1}^{t_\text{out}}\Big[\,
    \textstyle\sum_{d=1}^{d_i}\F_\text{step}\big(s_i^{(d)}\big)
    + (d_i{-}1)\,\F_\text{it}+\min(d_i,d_\text{max}{-}1)\,\F_\text{dec}
  \,\Big], \label{eq:tah-flops} \\[4pt]
  &\sum_{i=1}^{t_\text{out}}\Big[\,
    \textstyle\sum_{d=1}^{d_i}\M_\text{step}\big(s_i^{(d)}\big)
    + (d_i{-}1)\,\M_\text{it}+\min(d_i,d_\text{max}{-}1)\,\M_\text{dec}
  \,\Big], \label{eq:tah-mem}
\end{align}
where $s_i^{(d)}$ is the KV cache length visible at depth $d$ for the $i$-th output token under duo-causal attention (Equation~\ref{eq:duo_causal_attention}).
Since the KV cache at each depth is visible to all deeper levels, for $d_\text{max}{=}2$ we have
$s_i^{(1)} = t_\text{in} + i - 1$ (all prior tokens pass through depth~1) and
$s_i^{(2)} = s_i^{(1)} + |\{j < i : d_j {=} 2\}|$ (depth-1 entries plus prior depth-2 tokens).
Table~\ref{tab:tah-extra} lists the per-call cost of the inner transition and iteration decider; both are negligible relative to one backbone pass.
The overhead is therefore dominated by the average depth $\bar{d}=\frac{1}{t_\text{out}}\sum_i d_i$, giving approximately $\bar{d}\times$ the standard cost when $\bar{d}$ is small.

\xhdr{Extension to AlwaysThink} For AlwaysThink, every output token iterates to the maximum depth. When $d_i{=}d_\text{max}{=}2$ for all $i$, every depth-1 KV entry from output tokens has a corresponding depth-2 entry, giving $s_i^{(2)} \approx 2\,s_i^{(1)}$.
The two backbone passes incur $2\times$ the single-pass FFN cost, while their combined attention cost scales as $\mathcal{O}(s_i^{(1)}) + \mathcal{O}(s_i^{(2)}) \approx 3\,\mathcal{O}(s_i^{(1)})$, i.e., ${\sim}3\times$ the single-pass attention FLOPs.
Combining both components yields $2$-$3\times$ total overhead over Standard.


\begin{table}[ht]
\centering
\caption{Per-call cost of \name-specific components. The inner transition is applied $d_i{-}1$ times and the decider $\min(d_i,\,d_\text{max}{-}1)$ times per token. Both are negligible relative to one backbone pass.}
\label{tab:tah-extra}
\footnotesize
\setlength{\tabcolsep}{4pt}
\begin{tabular}{@{}llll@{}}
\toprule
\textbf{Component} & \textbf{Description} & \textbf{FLOPs} & \textbf{Memory (bytes)} \\
\midrule
Inner transition
  & $x_i^{(d+1)}{=}p_i^{(d)}E$ (Eq.~\ref{eq:ours_update})
  & $v + 2hk$
  & $(v {+} kh {+} h)\,b$ \\[3pt]
Iter.\ decider $\mathcal{I}_\phi$
  & MLP ($|\phi|$ params)
  & $2|\phi|$
  & $(|\phi| {+} \text{acts})\,b$ \\
\bottomrule
\end{tabular}
\end{table}

\xhdr{Results}
Tables~\ref{tab:token_iter_stats} and~\ref{tab:token_iter_stats_math} report the average number of input/output tokens and latent iterations per token across benchmarks.
We plug these statistics into the analysis to obtain the per-token costs in Table~\ref{tab:efficiency_analysis_theo}.
With only a small fraction of tokens iterating twice, \name incurs only 1.04-1.05$\times$ the cost of the Standard baseline.
In contrast, AlwaysThink sets $\bar{d}{=}2$, requiring 2.19-2.27$\times$ more computation and memory access.
These results confirm that \name achieves the performance benefits while avoiding the substantial efficiency penalty.


\begin{table}[t]
    \caption{Input tokens (shared across methods) and output token / iteration statistics for Standard, AlwaysThink, \name, and TaH+ (General-trained version).}
    \setlength\tabcolsep{4pt}
    \label{tab:token_iter_stats}
    \centering
    \begin{tabular}{lll|cc|cc|cc|cc|cc}
    \toprule
    & ~ & ~ & \multicolumn{2}{c|}{\textbf{Standard}} & \multicolumn{2}{c|}{\textbf{AlwaysThink}} &  \multicolumn{2}{c|}{\textbf{Ouro}}& \multicolumn{2}{c|}{\textbf{\name}} & \multicolumn{2}{c}{\textbf{TaH+}} \\
    \cmidrule(lr){4-5}\cmidrule(lr){6-7}\cmidrule(lr){8-9}\cmidrule(lr){10-11}\cmidrule(lr){12-13}
    \textbf{Param.} & \textbf{Dataset} & \textbf{In.} & Out. & Iter. & Out. & Iter.  & Out. &Iter.  & Out. & Iter. & Out. & Iter. \\
    \midrule
    \multirow{10}{*}{\textit{0.6B}}& AIME25        & 159& 7687& 1.00& 6885& 2.00& 7370&2.00& 7741& 1.04& 7677& 1.05\\
    & OlympiadBench & 100& 6937& 1.00& 6175& 2.00& 6733&2.00& 6926& 1.04& 6796& 1.04\\
    & AMC23         & 85& 6602& 1.00& 6332& 2.00& 6355&2.00& 6626& 1.04& 6688& 1.04\\
    & MATH500       & 71& 5364& 1.00& 5072& 2.00& 5274&2.00& 5189& 1.04& 5092& 1.04\\
    & GSM8K         & 61& 2144& 1.00& 1710& 2.00& 2039&2.00& 1895& 1.06& 1808& 1.06\\
    & GPQA          & 234& 6533& 1.00& 5998& 2.00& 6087&2.00& 6104& 1.11& 5966& 1.12\\
    & MMLU-STEM     & 97& 2302& 1.00& 2134& 2.00& 2201&2.00& 2057& 1.14& 2033& 1.14\\
    & HumanEval++   & 133& 4711& 1.00& 4411& 2.00& 4149&2.00& 4045& 1.05& 3621& 1.05\\
    & MBPP++        & 51& 4652& 1.00& 4964& 2.00& 4521&2.00& 3094& 1.05& 3175& 1.05\\
    \cmidrule(lr){2-13}
    % & Average       & 95& 5441& 1.00 & 5134& 2.00 & 5464& 1.06 & 5346& 1.06 \\
    & Average ratio & -- & 1.00$\times$ &1.00$\times$ & 0.93$\times$ & 2.00$\times$& 0.95$\times$& 2.00$\times$& 0.93$\times$& 1.06$\times$ & 0.91$\times$ &  1.07$\times$\\

    
    
    \midrule
    \multirow{10}{*}{\textit{1.7B}}& AIME25        & 159& 7436& 1.00& 7458& 2.00& 7621&2.00& 7799& 1.06& 7547& 1.07\\
    & OlympiadBench & 100& 6190& 1.00& 6318& 2.00& 6505&2.00& 6555& 1.06& 6403& 1.07\\
    & AMC23         & 85& 5918& 1.00& 5766& 2.00& 6261&2.00& 6423& 1.06& 5989& 1.07\\
    & MATH500       & 71& 4159& 1.00& 4347& 2.00& 4455&2.00& 4474& 1.06& 4076& 1.07\\
    & GSM8K         & 61& 1540& 1.00& 1775& 2.00& 1745&2.00& 1718& 1.07& 1597& 1.09\\
    & GPQA          & 234& 6269& 1.00& 6400& 2.00& 6380&2.00& 6729& 1.16& 6502& 1.18\\
    & MMLU-STEM     & 97& 1774& 1.00& 1729& 2.00& 2005&2.00& 1893& 1.14& 1769& 1.18\\
    & HumanEval++   & 133& 3875& 1.00& 3826& 2.00& 4177&2.00& 3905& 1.07& 3797& 1.09\\
    & MBPP++        & 51& 3438& 1.00& 3834& 2.00& 3682&2.00& 3181& 1.07& 3035& 1.09\\
    \cmidrule(lr){2-13}
    % & Average       & 95& 4868& 1.00 & 5452& 2.00 & 5192& 1.06 & 5116& 1.06\\
    & Average ratio & --& 1.00$\times$ & 1.00$\times$& 1.02$\times$& 2.00$\times$& 1.05$\times$ &2.00$\times$& 1.05$\times$ & 1.08$\times$ & 1.00$\times$ & 1.10$\times$ \\

    \midrule
    \multirow{10}{*}{\textit{4B}}& AIME25        & 159& 7007& 1.00& --& --& 7159&2.00& 7452& 1.04& 7282& 1.04\\
    & OlympiadBench & 100& 5699& 1.00& --& --& 5649&2.00& 6010& 1.05& 5898& 1.04\\
    & AMC23         & 85& 5155& 1.00& --& --& 5194&2.00& 5624& 1.04& 5638& 1.04\\
    & MATH500       & 71& 3426& 1.00& --& --& 3434&2.00& 3657& 1.05& 3696& 1.04\\
    & GSM8K         & 61& 1345& 1.00& --& --& 1306&2.00& 1564& 1.06& 1498& 1.06\\
    & GPQA          & 234& 6200& 1.00& --& --& 6133&2.00& 6106& 1.12& 6120& 1.12\\
    & MMLU-STEM     & 97& 1589& 1.00& --& --& 1541&2.00& 1632& 1.11& 1577& 1.11\\
    & HumanEval++   & 133& 3047& 1.00& --& --& 2776&2.00& 3312& 1.05& 3030& 1.09\\
    & MBPP++        & 51& 2673& 1.00& --& --& 2402&2.00& 2633& 1.05& 2584& 1.05\\
    \cmidrule(lr){2-13}
    % & Average       & 95& 4868& 1.00 & 5452& 2.00 & 5192& 1.06 & 5116& 1.06\\
   & Average ratio & --& 1.00$\times$& 1.00$\times$& --& -- & 0.98$\times$&2.00$\times$& 1.05$\times$ &1.06$\times$ & 1.03$\times$ & 1.07$\times$\\


    \bottomrule
    \end{tabular}
\end{table}

\begin{table}[t]
    \caption{Input tokens (shared across methods) and output token / iteration statistics for Standard, AlwaysThink, \name, and TaH+ (Math-trained version).}
    \setlength\tabcolsep{4pt}
    \label{tab:token_iter_stats_math}
    \centering
    \begin{tabular}{lll|cc|cc|cc|cc}
    \toprule
    & ~ & ~ & \multicolumn{2}{c|}{\textbf{Standard}} & \multicolumn{2}{c|}{\textbf{AlwaysThink}} & \multicolumn{2}{c|}{\textbf{\name}} & \multicolumn{2}{c}{\textbf{TaH+}} \\
    \cmidrule(lr){4-5}\cmidrule(lr){6-7}\cmidrule(lr){8-9}\cmidrule(lr){10-11}
    \textbf{Param.} & \textbf{Dataset} & \textbf{In.} & Out. & Iter. & Out. & Iter. & Out. & Iter. & Out. & Iter. \\
    \midrule
    \multirow{6}{*}{\textit{0.6B}}& AIME25        & 159& 7450 & 1.00 & 7316 & 2.00 & 7648 & 1.05 & 7486 & 1.06 \\
    & OlympiadBench & 100& 6599 & 1.00 & 6622 & 2.00 & 6631 & 1.09 & 6513 & 1.06 \\
    & AMC23         & 85& 6377 & 1.00 & 6368 & 2.00 & 6242 & 1.05 & 6145 & 1.05 \\
    & MATH500       & 71& 4823 & 1.00 & 5350 & 2.00 & 4877 & 1.05 & 4793 & 1.06 \\
    & GSM8K         & 61& 1955 & 1.00 & 2844 & 2.00 & 1923 & 1.07 & 1791 & 1.07 \\
    \cmidrule(lr){2-11}
    % & Average       & 95& 5441& 1.00 & 5134& 2.00 & 5464& 1.06 & 5346& 1.06 \\
    & Average ratio & --& 1.00$\times$& 1.00$\times$& 1.02$\times$& 2.00$\times$& 1.00$\times$& 1.06$\times$& 0.97$\times$& 1.06$\times$\\


    \midrule
    \multirow{6}{*}{\textit{1.7B}}& AIME25        & 159& 7195 & 1.00 & 7173 & 2.00 & 7496 & 1.06 & 7498 & 1.06\\
    & OlympiadBench & 100& 6008 & 1.00 & 6484 & 2.00 & 6387 & 1.06 & 6258 & 1.06 \\
    & AMC23         & 85& 5681 & 1.00 & 7543 & 2.00 & 6122 & 1.04 & 5852 & 1.06 \\
    & MATH500       & 71& 4004 & 1.00 & 4414 & 2.00 & 4233 & 1.06 & 4286 & 1.06 \\
    & GSM8K         & 61& 1451 & 1.00 & 1644 & 2.00 & 1721 & 1.08 & 1686 & 1.08 \\
    \cmidrule(lr){2-11}
    % & Average       & 95& 4868& 1.00 & 5452& 2.00 & 5192& 1.06 & 5116& 1.06\\
    & Average ratio & --& 1.00$\times$& 1.00$\times$& 1.13$\times$& 2.00$\times$& 1.09$\times$& 1.06$\times$ & 1.07$\times$& 1.06$\times$ \\

    \bottomrule
    \end{tabular}
\end{table}

\begin{table}[t]
    \caption{Decoding computation (GFLOPs) and memory access (GB) per output token for Standard, AlwaysThink, \name, and TaH+.}
    \setlength\tabcolsep{3pt}
    \label{tab:efficiency_analysis_theo}
    \centering
    \begin{tabular}{ll|cc|cc|cc|cc}
    \toprule
    & & \multicolumn{2}{c|}{\textbf{Standard}} & \multicolumn{2}{c|}{\textbf{AlwaysThink}} & \multicolumn{2}{c|}{\textbf{\name}} & \multicolumn{2}{c}{\textbf{TaH+}} \\
    \cmidrule(lr){3-4}\cmidrule(lr){5-6}\cmidrule(lr){7-8}\cmidrule(lr){9-10}
    \textbf{Param.} & \textbf{Dataset} & Comp.& Mem.& Comp.& Mem.& Comp.& Mem.& Comp.& Mem.\\

    \midrule
    \multirow{6}{*}{\textit{0.6B}}& AIME25        & 1.47 & 1.38 & 3.35 & 3.14& 1.52 & 1.43& 1.57 & 1.47 \\
    & OlympiadBench & 1.41 & 1.32 & 3.21 & 3.02& 1.51 & 1.42& 1.50 & 1.41 \\
    & AMC23 & 1.40 & 1.31 & 3.17 & 2.97& 1.43 & 1.34& 1.46 & 1.37 \\
    & MATH500 & 1.31 & 1.22 & 2.98 & 2.80& 1.35 & 1.26& 1.39 & 1.31 \\
    & GSM8K & 1.14 & 1.06 & 2.54 & 2.38& 1.19 & 1.12& 1.22 & 1.14 \\
    \cmidrule(lr){2-10}
    & Average ratio & 1.00$\times$ & 1.00$\times$ & 2.27$\times$ & 2.27$\times$& 1.04$\times$ & 1.04$\times$& 1.06$\times$ & 1.06$\times$ \\
    \midrule
    \multirow{6}{*}{\textit{1.7B}}& AIME25        & 4.31 & 4.03 & 9.45 & 8.83& 4.51 & 4.21& 4.64 & 4.34 \\
    & OlympiadBench & 4.16 & 3.88 & 9.18 & 8.58& 4.36 & 4.07& 4.48 & 4.18 \\
    & AMC23 & 4.12 & 3.85 & 9.54 & 8.91& 4.24 & 3.96& 4.43 & 4.13 \\
    & MATH500 & 3.92 & 3.66 & 8.45 & 7.89& 4.10 & 3.83& 4.23 & 3.95 \\
    & GSM8K & 3.62 & 3.38 & 7.48 & 6.98& 3.87 & 3.61& 3.98 & 3.72 \\
    \cmidrule(lr){2-10}
    & Average ratio & 1.00$\times$ & 1.00$\times$ & 2.19$\times$ & 2.19$\times$& 1.05$\times$ & 1.05$\times$& 1.08$\times$ & 1.08$\times$ \\
    \bottomrule
    \end{tabular}
\end{table}

%%%%%%%%%%%%%%%%% alwaysthink: FREE THIS %%%%%%%%%%%%%%%%

% \subsubsection{Attention Ablation for AlwaysThink}
% \label{sec:appendix/ablation_alwaysthink_attention}

% \begin{table}[tb]
%     \caption{Ablation of attention mechanisms under the AlwaysThink policy (TaH-0.6B).}
%     \label{tab:ablation_arch_alwaysthink}
%     \centering
%     \setlength\tabcolsep{6pt}
%     \begin{tabular}{cccc|cccc}
%         \toprule
%         \textbf{Iter. Policy} & \textbf{Attention} & \textbf{LoRA} & \textbf{Residual} & \textbf{MATH500} & \textbf{AMC23} & \textbf{Olym.} & \textbf{Average} \\
%         \midrule
%         Iter. Decider & Duo-causal & \cmark & \cmark & \textbf{51.2} & \textbf{32.5} & \textbf{23.9} & \textbf{35.9}{\scriptsize$/+$0.0} \\
%         \midrule
%         AlwaysThink & Duo-causal & \cmark & \cmark & 32.8 & 15.6 & 10.2 & 19.5{\scriptsize$/-$16.4} \\
%         AlwaysThink & Causal     & \cmark & \cmark & 44.6 & 20.6 & 16.0 & 27.1{\scriptsize$/-$8.8} \\
%         \bottomrule
%     \end{tabular}
% \end{table}

% Duo-causal attention facilitates cross-iteration information flow, which is critical for dynamic depths. However, in fixed-depth settings (AlwaysThink) where full context is already available within the current iteration, extending the context via duo-causal attention can introduce noise and degrade performance. To provide insights for future research on the trade-offs between connectivity and noise, we compare attention mechanisms under the AlwaysThink policy in Table~\ref{tab:ablation_arch_alwaysthink}.

\subsection{Additional Analysis}
\label{sec:appendix/exp/analysis}

\subsubsection{Training Scheme Dynamics}
\label{sec:appendix/exp/training_dynamics}
Figure~\ref{fig:training_scheme_dynamics} complements the training-scheme ablations in Table~\ref{tab:ablation} with validation-perplexity curves.
Token-only supervision with a fixed oracle policy trains stably, while token+latent supervision and decider-based training lag behind.
The dynamic oracle can reduce validation perplexity but is unstable during generation, matching its poor downstream accuracy.
For the oracle discrepancy metric, top-1 mismatch gives the best trajectory under matched continuation budgets, supporting our default policy. Note that these training-scheme dynamics reflect a performance ceiling rather than the actual downstream task performance.

\begin{figure}[tb]
    \centering
    \begin{subfigure}[t]{0.48\textwidth}
        \centering
        \includegraphics[width=\textwidth]{figure/eval_loss_ablation.pdf}
        \caption{Supervision type and iteration policy.}
        \label{fig:ablation_validation_ppl}
    \end{subfigure}
    \hfill
    \begin{subfigure}[t]{0.48\textwidth}
        \centering
        \includegraphics[width=\textwidth]{figure/eval_loss_label_ablation.pdf}
        \caption{Oracle discrepancy metric.}
        \label{fig:loss_difficulty_metric}
    \end{subfigure}
    \caption{Validation-perplexity dynamics for the training-scheme ablations in Table~\ref{tab:ablation}. (a) Supervision and iteration-policy choices. (b) Oracle discrepancy metrics under matched continuation budgets.}
    \label{fig:training_scheme_dynamics}
\end{figure}


\subsubsection{Iteration Label Analysis}
\label{sec:appendix/exp/oracle_label}

\xhdr{Iteration label identifiability}
We first examine whether the oracle iteration labels expose stable token-level signals.
We compute the token entropy of continue-labeled and stop-labeled tokens across three diverse subsets of the Open-R1 dataset (Math, Science, and Code).
As shown in Figure~\ref{fig:label_hard_tokens}, continue-labeled tokens exhibit a universal signature of significantly higher entropy ($>5\times$) than stop-labeled tokens.
This distinct separation confirms that iteration need is an intrinsic, robustly identifiable property of the model's predictive state, rather than a complex, task-specific pattern.
Given this clear signal, the neural iteration decider can learn reliable classification strategies that generalize well across different domains.

\xhdr{Cross-model label consistency}
We next analyze whether iteration labels are stable across reference model scales.
Figure~\ref{fig:venn_overlap} shows substantial overlap; for example, the 1.7B reference identifies 81\% of the iteration-selected tokens found by the 4B reference.
To assess the quality of this overlap, we compare the cross-entropy of overlap and non-overlap iteration-selected tokens in Figure~\ref{fig:ref_overlap_ce}.
Overlap tokens have consistently higher loss ($\approx 2.0\times$), indicating a stable core of genuinely difficult token positions.
These results show that iteration labels expose stable uncertainty signals, motivating the generalization analysis of the learned iteration decider below.

\begin{figure}[tb]
    \centering
    \hfill
    \begin{subfigure}[t]{0.48\textwidth}
        \centering
        \includegraphics[width=0.5\textwidth]{figure/venn_0.6_1.7.pdf}
        \caption{Overlap between Qwen3-0.6B and 1.7B}
        \label{fig:venn_0_6_1_7}
    \end{subfigure}
    \hfill
    \begin{subfigure}[t]{0.48\textwidth}
        \centering
        \includegraphics[width=0.5\textwidth]{figure/venn_1.7_4.pdf}
        \caption{Overlap between Qwen3-1.7B and 4B}
        \label{fig:venn_1_7_4}
    \end{subfigure}
    \hfill
    \caption{Venn diagrams illustrating the overlap of iteration-selected tokens between different reference models. The high overlap proportions indicate that iteration labels are largely consistent across model scales.}
    \label{fig:venn_overlap}
\end{figure}

\begin{figure}[tb]
    \centering
    \begin{subfigure}[t]{0.48\textwidth}
        \centering
        \includegraphics[width=0.7\textwidth]{figure/experiment/ce_1.7_vs_0.6.pdf}
        \caption{Qwen3-1.7B: cross-entropy of overlap vs. non-overlap iteration-selected tokens (w.r.t. Qwen3-0.6B).}
        \label{fig:ce_1_7_vs_0_6}
    \end{subfigure}
    \hfill
    \begin{subfigure}[t]{0.48\textwidth}
        \centering
        \includegraphics[width=0.7\textwidth]{figure/experiment/ce_0.6_vs_1.7.pdf}
        \caption{Qwen3-0.6B: cross-entropy of overlap vs. non-overlap iteration-selected tokens (w.r.t. Qwen3-1.7B).}
        \label{fig:ce_0_6_vs_1_7}
    \end{subfigure}

    \vspace{0.4em}

    \begin{subfigure}[t]{0.48\textwidth}
        \centering
        \includegraphics[width=0.7\textwidth]{figure/experiment/ce_4_vs_1.7.pdf}
        \caption{Qwen3-4B: cross-entropy of overlap vs. non-overlap iteration-selected tokens (w.r.t. Qwen3-1.7B).}
        \label{fig:ce_4_vs_1_7}
    \end{subfigure}
    \hfill
    \begin{subfigure}[t]{0.48\textwidth}
        \centering
        \includegraphics[width=0.7\textwidth]{figure/experiment/ce_1.7_vs_4.pdf}
        \caption{Qwen3-1.7B: cross-entropy of overlap vs. non-overlap iteration-selected tokens (w.r.t. Qwen3-4B).}
        \label{fig:ce_1_7_vs_4}
    \end{subfigure}

    \caption{Token-level cross-entropy distributions of overlap and non-overlap iteration-selected tokens across different reference model pairs. For each pair of reference models (e.g., Qwen3-1.7B and Qwen3-0.6B), we plot the cross-entropy of tokens selected for iteration by both models (overlap) and by only one model (non-overlap) on both reference models.}
    \label{fig:ref_overlap_ce}
\end{figure}


\begin{figure}[tb]
    \centering
    \begin{subfigure}[t]{0.32\textwidth}
        \centering
        \includegraphics[width=\textwidth]{figure/experiment/label_math.pdf}
        \caption{Math}
        \label{fig:label_math}
    \end{subfigure}
    \hfill
    \begin{subfigure}[t]{0.32\textwidth}
        \centering
        \includegraphics[width=\textwidth]{figure/experiment/label_science.pdf}
        \caption{Science}
        \label{fig:label_science}
    \end{subfigure}
    \hfill
    \begin{subfigure}[t]{0.32\textwidth}
        \centering
        \includegraphics[width=\textwidth]{figure/experiment/label_code.pdf}
        \caption{Code}
        \label{fig:label_code}
    \end{subfigure}
    \caption{Output-logit entropy distribution at the first iteration of \name, categorized by oracle iteration labels on the Open-R1 validation set (Math, Science, Code). The distinct separation between distributions confirms that \name's internal logits provide a strong, task-agnostic signal for identifying tokens that require additional iteration.}
    % , explaining the high generalizability of the iteration decider.}
    \label{fig:label_hard_tokens}
\end{figure}


\subsubsection{Iteration Decider Analysis}

\xhdr{Cross-domain behavior}
We evaluate the iteration decider, trained on the general Open-R1 corpus, across three validation subsets (Math, Code, and QA) to quantify its robustness and cross-domain generalizability.
As summarized in Table~\ref{tab:decider_generalization}, the decider maintains high decision accuracy across all domains without any retraining.

Despite being invoked on only 7.8-26.6\% of tokens, the decider consistently yields 5.2-9.0\% absolute accuracy gains over the standard single-pass baseline on all three domains.
Moreover, the decider automatically adjusts its iteration rate according to task difficulty: it iterates more frequently on QA (26.6\%) than on Math (7.8\%), even under a fixed threshold $c_{\text{threshold}}=0.9$.
This behavior indicates that the decider responds to intrinsic uncertainty signals in the model's predictive distribution rather than memorizing domain-specific patterns.

\xhdr{Iteration decision error}
We further analyze how specific decision mistakes affect end-to-end response quality.
Because the learned decider is imperfect, as shown in Figure~\ref{fig:training_decider}, we randomly inject errors into the oracle iteration-decider predictions at different rates.
Formally, we denote the original oracle prediction as the \textit{label} $l \in \{0,1\}$ and the altered prediction as the \textit{output} $o \in \{0,1\}$.
We define the \textit{iter. error} as the total proportion of deliberately introduced errors:
\begin{equation}
    \text{iter. error} = P(l \neq o) =
    \underbrace{P(l=1,o=0)}_{\textit{underthink rate}}
    + \underbrace{P(l=0,o=1)}_{\textit{overthink rate}}.
\end{equation}
We further distinguish the impacts of overthinking and underthinking.
Here, overthinking refers to cases where the decider incorrectly signals \textit{continue}, while underthinking corresponds to cases where it incorrectly signals \textit{stop}.
Table~\ref{tab:oracle_with_noise} shows how \name's MATH100 accuracy varies with different iteration error rates.
We quantify these effects by fitting a linear model to the data:
$$
\text{accuracy} = -1.41 \times \text{underthink rate} - 2.73 \times \text{overthink rate} + 0.81.
$$
This analysis indicates that inaccurate iteration decisions are the main factor behind the performance gap between \name and its oracle variant, with overthinking being the dominant source of performance gaps.

\begin{table}[tb]
    \centering
    \begin{minipage}[t]{0.47\textwidth}
        \vspace{0pt}
        \centering
        \caption{Iteration decider behavior and downstream gains on different validation subsets. The decider is trained once on general Open-R1 and evaluated without retraining.}
        \label{tab:decider_generalization}
        \footnotesize
        \setlength{\tabcolsep}{3pt}
        \begin{tabular}{lccc}
        \toprule
        \textbf{Metric} & \textbf{Math} & \textbf{Code} & \textbf{QA} \\
        \midrule
        Iteration Percentage & 7.8\% & 10.7\% & 26.6\% \\
        Iteration Accuracy & 86.7\% & 82.3\% & 76.6\% \\
        \makecell[l]{Benchmark Gain\\over Standard} & \textbf{+5.2\%} & \textbf{+9.0\%} & \textbf{+5.8\%} \\
        \bottomrule
        \end{tabular}
    \end{minipage}
    \hfill
    \begin{minipage}[t]{0.50\textwidth}
        \vspace{0pt}
        \centering
        \caption{\name performance under different iteration-decider error rates. All values are reported in percentages.}
        \label{tab:oracle_with_noise}
        \footnotesize
        \setlength{\tabcolsep}{3pt}
        \begin{tabular}{ccc|c}
        \toprule
        \makecell{\textbf{Iter.}\\\textbf{Error (\%)}} & \makecell{\textbf{Underthink}\\\textbf{(\%)}} & \makecell{\textbf{Overthink}\\\textbf{(\%)}} & \makecell{\textbf{MATH100}\\\textbf{Acc. (\%)}}\\
        \midrule
        0.0 & 0.0 & 0.0 & 80.0\\
        2.8 & 2.8 & 0.0 & 78.0\\
        10.0 & 1.5 & 8.5 & 55.4\\
        15.0 & 2.1 & 12.9 & 45.2\\
        20.0 & 2.5 & 17.5 & 27.1\\
        22.1 & 0.0 & 22.1 & 21.6\\
        \bottomrule
        \end{tabular}
    \end{minipage}
\end{table}


\subsubsection{Token Alternation Pattern}
\label{sec:appendix/exp/flow}

We analyze tokens that most frequently trigger a second iteration ("think-twice" tokens). For each token type $t$, we compute the continuation rate
\[
\Pr\big(c_i^{(1)} > c_{\text{threshold}} \mid t_i = t\big),
\]
using the inference threshold $c_{\text{threshold}}=0.9$ (Section~\ref{sec:training}). We estimate this quantity on the Open-R1 validation set and, for diagnostics, randomly sample 10K token positions (\(\approx\)0.4\% of tokens) to track whether the next-token prediction switches between depth 1 and depth 2. This setting quantifies which token types most often trigger an additional iteration and how often iteration alters the predicted next token.


\begin{table}[tb]
    \centering
    \caption{Conditional probabilities of continuation confidence and next-token distribution.}
    \begin{tabular}{lclc}
    \toprule
    \textbf{Token $T_1$} & $P(c^{(1)} > c_{\text{threshold}} \mid t^{(1)}=T_1)$ & \textbf{Token $T_2$} & $P(t^{(2)}=T_2 \mid t^{(1)}=T_1)$ \\
    \midrule
    \multirow{3}{*}{But} & \multirow{3}{*}{34.3\%} & So        & 13.63\% \\
                         &                          & Wait      & 12.17\% \\
                         &                          & Therefore & 8.95\%  \\
    \midrule
    \multirow{3}{*}{So}  & \multirow{3}{*}{17.7\%} & So        & 28.17\% \\
                         &                          & Therefore & 13.67\% \\
                         &                          & But       & 4.89\%  \\
    \bottomrule
    \end{tabular}
    \label{tab:token_probs}
\end{table}



\subsubsection{Duo-causal Attention Pattern}
\label{sec:appendix/exp/attention}

We perform forward computation on 100 samples, each with a length of 128 tokens, and visualize the learned attention patterns of the \name model during the second iteration.

\xhdr{Qualitative analysis}
Figure~\ref{fig:attention} shows the average attention weights of three representative heads.
The left panel illustrates a head that mainly attends to first-iteration keys; the middle panel shows one focusing on second-iteration keys; and the right panel displays a head with balanced attention across both iterations.
These examples demonstrate that the duo-causal attention mechanism enables the model to automatically learn diverse cross-depth attention strategies.

\begin{figure}[tb]
    \centering
    \includegraphics[width=\textwidth]{figure/experiment/attention.pdf}
    \caption{\name duo-causal attention pattern.}
    \label{fig:attention}
\end{figure}

\xhdr{Quantitative analysis}
Figure~\ref{fig:K1_attention_layer_curve} further quantifies, for each layer, how much attention mass second-iteration queries allocate to first-iteration keys.
Following the \name framework, let $a_{\ell,h}\!\big((i,d_q)\!\to\!(j,d_{kv})\big)$ denote the attention weight at layer $\ell$, head $h$, from the query at position $i$, depth $d_q$ to the key at position $j$, depth $d_{kv}$.
As an illustrative example, consider query $(3,2)$ highlighted in Figure~\ref{fig:duo_attention_expand}(c): it attends to three depth-1 keys and two depth-2 keys. We want to quantify how attention mass is distributed between these two groups.
We define the \emph{cross-iteration attention mass} for a single query position as
\[
m_{i,\ell,h}^{(2\to 1)} \;=\; \sum_{j \le i} a_{\ell,h}\!\big((i,2)\!\to\!(j,1)\big),
\]
and average over all depth-2 query positions within each sequence $s$:
\[
\bar{m}_{s,\ell,h}^{(2\to 1)}
\;=\;
\frac{1}{\lvert Q_s^{(2)} \rvert}
\sum_{i \in Q_s^{(2)}} m_{i,\ell,h}^{(2\to 1)},
\]
where $Q_s^{(2)}$ is the set of token positions that iterate to depth $2$ in sequence $s$.
The curve in Figure~\ref{fig:K1_attention_layer_curve} plots, for each layer, the mean of $\bar{m}_{s,\ell,h}^{(2\to 1)}$ over all sequence-head pairs $(s,h)$, with shading indicating one standard deviation.

\xhdr{Findings}
A higher $\bar{m}^{(2\to 1)}$ indicates that depth-2 queries rely more on first-iteration keys, whereas a lower value indicates greater reliance on same-iteration keys.
The results reveal substantial layer-wise heterogeneity: lower layers place relatively less mass on first-iteration keys, while deeper layers show higher and varied cross-depth reuse.

\begin{figure}[tb]
    \centering
    \includegraphics[width=0.55\textwidth]{figure/experiment/k1_sum_attention_analysis.pdf}
    \caption{\name layer-wise cross-iteration attention mass $\bar{m}^{(2\to 1)}$ (mean $\pm$ one standard deviation over sequences and heads).}
    \label{fig:K1_attention_layer_curve}
\end{figure}

%%%%%%%%%%%% implementation %%%%%%%%%%%%
\subsection{Implementation Details}

%%%%%%%%%%%%%%%%% attention %%%%%%%%%%%%%%%%%
\subsubsection{Duo-causal attention implementation}
\label{sec:appendix/duo-causal-attention}

Figure~\ref{fig:duo_attention_expand} illustrates the implementation of duo-causal attention, with the formal definitions provided below.

\textbf{(1) KV cache concatenation.} At depth $d$, we form the visible K/V sequence by concatenating all shallower-to-current depths along the sequence dimension:
\[ \mathrm{KV}^{(\le d)} \,=\, [\, \mathrm{KV}^{(1)}\,;\,\mathrm{KV}^{(2)}\,;\,\cdots\,;\,\mathrm{KV}^{(d)} \,]. \]
This realizes the accessible set in Equation~\ref{eq:duo_causal_attention}, allowing deeper iterations to access all shallower iterations while preserving positional causality. 
The KV cache is managed by iteration depth during decoding, as shown in Figure~\ref{fig:duo_attention_expand}(b). The fragmented KV-cache management strategy is standard in existing LLM serving systems~\citep{kwon2023vllm, zheng2024sglang}.

\textbf{(2) Two-dimensional causal mask.} For a query $(i,d)$, a key $(j,k)$ is attendable iff $j\le i$ and $k\le d$. We implement this as an additive attention mask with $0$ for allowed entries and $-\infty$ otherwise, enforcing positional and iteration causality jointly. 
Figure~\ref{fig:duo_attention_expand}(c) visualizes the landscape of the duo-causal attention mask.
When $d=1$ for all tokens, the rule reduces to standard causal attention.

\textbf{(3) Compatibility with efficient attention.} The mask is provided in the standard additive form and the concatenated K/V tensors remain contiguous along the sequence dimension, matching the usual scaled dot-product attention interface.
As a result, duo-causal attention is directly compatible with optimized kernels such as FlashAttention, without kernel modifications.

%%%%%%%%%%%%%%%%% related work %%%%%%%%%%%%%%%%%
\subsection{Additional Related Work}
Instead of using the shared model parameters multiple times through latent iteration, previous work also proposes layer skipping methods for dynamic compute allocation.

\xhdr{Layer skipping}
Layer skipping aims to accelerate LLM inference by dynamically bypassing certain layers for specific tokens.
Some methods use a learnable module to make real-time skipping decisions. MoD~\citep{raposo2024mixture} uses a top-k router to select a subset of tokens for processing, while FlexiDepth~\citep{luo2025adaptive} uses a plug-in router to determine whether a layer should be bypassed. Others use a fixed strategy to skip layers. SkipDecode~\citep{del2023skipdecode} enforces a monotonically decreasing number of active layers during generation.
However, these methods still require loading the entire model's parameters, resulting in a large memory access overhead.
Instead of skipping some layers, \name adds computational depth by allowing core tokens to undergo multiple refinement iterations.
This approach provides greater computational depth without increasing the model's parameter count.

% \xhdr{Concurrent works}
% \todo{explain concurrent works, including MoR, xxx}
% Notably, previous work~\citep{mohtashami2023cotformer} also explores dynamic recursion. But it mainly perceives the problem by letting the sequence length grow with the number of repeats, while our approach focuses on KV-Cache management and attention pattern manipulation, instead of expanding the entire sequence.

%%%%%%%%%%%% Limitations and future work %%%%%%%%%%%%
\subsection{Limitations and Future Work}
\xhdr{Comparison with official Qwen3 models}
Official Qwen3 models are trained on different data distributions and scales, and use different training procedures, including on-policy distillation~\citep{yang2025qwen3}. By contrast, our models use SFT only on limited, publicly accessible data. Consequently, performance may differ between the two.

\xhdr{Future work}
We consider advanced training techniques largely orthogonal to \name. Promising directions include (1) integrating \name with online distillation to provide stronger supervision; (2) applying reinforcement learning to optimize the compute–accuracy trade-off and to learn continuation policies beyond imitation. Due to resource constraints, this work focuses on SFT with public data; we leave these combinations for future work.


%%%%%%%%%%%%%%%%% LLM %%%%%%%%%%%%%%%%%
\subsection{Use of Large Language Models (LLMs)}
Large language models (e.g., GPT-5, Gemini) were used exclusively for proofreading and minor copyediting. They were not employed to generate scientific content, contribute to research ideation, or design methods. All research ideas, methods, and conclusions are solely the responsibility of the authors.

%%%%%%%%%%%%%%%%%%%% case study %%%%%%%%%%%%%%%%%%%%%%%%%%
\subsection{Examples}

We use an example question from AIME25 and responses from Standard, TaH and AlwaysThink to provide an intuitive understanding of our method. Only TaH can successfully answer this question. Blue-highlighted tokens indicate positions that undergo two iterations.

%%%%%%%%%%%%%%%%% Example %%%%%%%%%%%%%%%%%
\begin{bluebox}[Question]{width=\linewidth, valign=top}
In $\triangle ABC$, points $A,D,E$, and $B$ lie in that order on side $\overline{AB}$ with $AD=4, DE=16$, and $EB=8$. Points $A,F,G$, and $C$ lie in that order on side $\overline{AC}$ with $AF=13, FG=52$, and $GC=26$. Let $M$ be the reflection of $D$ through $F$, and let $N$ be the reflection of $G$ through $E$. Quadrilateral $DEGF$ has area 288. Find the area of the heptagon $AFNBCEM$.
\end{bluebox}

\clearpage

\begin{figure*}[t]
\centering
\begin{minipage}[t]{0.32\textwidth}
\begin{bluebox}[Standard Output]{width=\linewidth, height=0.95\textheight, valign=top}

Okay, let me try to figure out this problem step by step. It's about a heptagon formed by reflecting points and finding its area. \ldots
We have triangle $ABC$. On side $AB$, points $A, D, E$, and $B$ are in that order. The lengths are given as $AD = 4$, $DE = 16$, and $EB = 8$. So $AB = 28$ \ldots
On side $AC$, points $A, F, G$, and $C$ are in that order with $AF = 13$, $FG = 52$, and $GC = 26$. So $AC = 91$ \ldots
Let me place point $A$ at $(0,0)$ and $B$ at $(28,0)$ \ldots Let $C$ be at $(91, h)$ \ldots Then
$D=(4,0)$, $E=(20,0)$, $F=\big(13,\frac{13h}{91}\big)$, $G=\big(65,\frac{65h}{91}\big)$ \ldots

Area of quadrilateral $DEGF$ can be calculated \ldots Applying the shoelace formula \ldots
Area $= \frac{1}{2}\cdot \frac{1248h}{91} = \frac{624h}{91} \;\mathbf{=}\; 6h$

So with $288$:
$6h = 288 \Rightarrow h = \frac{288}{6} = 48$

Reflections: $M$ is the reflection of $D$ through $F$ and $N$ is the reflection of $G$ through $E$ \ldots
$M=(22,0)$, $N=(-39,0)$, $B=(28,0)$, $C=(91,48)$, $E=(20,0)$ \ldots

Using the shoelace formula for $AFNBCEM$ in order \ldots
Sum $= 0 + 0 + 0 + 1344 + (-960) + 0 + 0 = 384$

Take the absolute value and multiply by $1/2$: Area $= \frac{1}{2}\cdot |384| = 192$

\ldots

\textbf{Final Answer} $\boxed{192}$ \xmark

\end{bluebox}
\end{minipage}\hfill
\begin{minipage}[t]{0.32\textwidth}
\begin{bluebox}[TaH Output]{width=\linewidth,  
height=0.95\textheight,  valign=top}


Okay\textcolor{blue}{,} let's try to tackle this geometry problem. It seems a bit \textcolor{blue}{ complex} with all these points and reflections\textcolor{blue}{,} \ldots
First, let me visualize triangle ABC. Points D, E, and B are on side AB with AD = 4, DE = 16, and EB = 8. \ldots
Let me set up a coordinate system\textcolor{blue}{. Let}'s place point A at the origin $(0, 0)$ \ldots Then, since $AB$ is a line segment, let \textcolor{blue}{ me} assume $AB$ is along the $x$-axis. So point $B$ would be at $(28, 0)$ \ldots
Now\textcolor{blue}{,} points $F$\textcolor{blue}{, G}, $C$ are on $AC$ \ldots Let \textcolor{blue}{me} define coordinates for $C$ as $(c_x, c_y)$ \ldots

\ldots

Compute the \textcolor{blue}{sum} of $x_i y_{i+1} - x_{i+1} y_i$ for each $i$ \ldots

Sum of all terms: $0 + \frac{1300 c_y}{91} + 0  -> \frac{52 c_y}{91} = \frac{(1300 - 52)c_y}{91} = \frac{1248 c_y}{91}$

Then, the area is $\frac{1}{2} \cdot \left|\frac{1248 c_y}{91}\right| = \frac{624 c_y}{91}$

\ldots

Compute the area of this heptagon \ldots 

Using the \textcolor{blue}{shoelace} formula \ldots

Term2: $-\!240$ \quad Term3: $840$ \quad Term4: $1176$ \quad Term5: $-\!840$ \quad Term6: $240$ \ldots

Sum: $0 - 240 + 840 + 1176 - 840 + 240 + 0 = 1176$

Take the absolute value and multiply by $1/2$:
Area $= \frac{1}{2}\cdot |1176| = 588$

\ldots

\textbf{Final Answer} $\boxed{588}$ \cmark

\end{bluebox}
\end{minipage}\hfill
\begin{minipage}[t]{0.32\textwidth}
\begin{bluebox}[AlwaysThink Output]{width=\linewidth, 
height=0.95\textheight, valign=top}

\begingroup
\color{blue}
Okay, so I need to find the area of the heptagon AFNBCEM. Let me first try to visualize the problem. There's triangle ABC. On side AB, there are points A, D, E, and B. The order is A, D, E, B. The distances given are AD=4, DE=16, and EB=8. So, starting at A, moving 4 units to D, then 16 units to E, then 8 units to B. So, the total length of AB is AD + DE + EB = 4 + 16 + 8 = 28 units. Got that.

\ldots

Points in order: D(4,0), E(20,0), F(13,h), G(65,h), back to D(4,0).

Compute the terms:

\ldots

Sum these terms: $0 + 20h - 52h - 4h = -36h$

Take the absolute value and multiply by 1/2: $Area = (1/2)*| -36h | = (1/2)*36|h| = 18|h|$

But the area is given as 288, so $18|h| = 288 \Rightarrow |h| = 16$.

\ldots

Now, apply the shoelace formula \ldots
Therefore, total sum is $8640/17$

Take the absolute value and multiply by 1/2: $Area = (1/2)*|8640/17| = 8640/34 = 4320/17 \approx 254.1176$

\ldots

Therefore, the answer is $4320/17$.

\endgroup
\textbf{Final Answer} $\boxed{\frac{4320}{117}}$ \xmark
\end{bluebox}
\end{minipage}
\end{figure*}












% \begin{figure}[h]
%     \begin{minipage}{0.68\textwidth}
%         \centering
%         \includegraphics[width=\textwidth]{figure/idea.png}
%         \caption{Like verbal reasoning, fixed-depth latent reasoning may overthink simple tokens or underthink hard ones. \name thinks deeply only on hard tokens to optimize overall reasoning quality.}
%         \label{fig:idea}
%     \end{minipage}\hfill
%     \begin{minipage}{0.30\textwidth}
%         \centering
%         \label{tab:oracle}
%         \begin{tabular}{l|c}
%         \toprule
%         \textbf{Method} & \textbf{Acc.} \\
%         \midrule
%         Direct & 52\% \\
%         Think-twice &  38\%  \\
%         Think at hard & \textbf{80\%}  \\
%         \bottomrule
%         \end{tabular}
%         \captionof{table}{
%             math500 100 sample performance
%         }
%         \label{tab:think_oracle}
%     \end{minipage}
% \end{figure}

% Table
% Limit response length to 4K
% Always iter@2 training
% Compare the output from reference model and current output
% Teacher: DeepSeek-R1-7B
